\documentclass[12pt, a4paper]{article}

\usepackage[utf8]{inputenc}
\usepackage[spanish]{babel}
\usepackage{amsmath}
\usepackage{amssymb}
\usepackage{booktabs}
\usepackage{float}
\usepackage{geometry}
\usepackage{hyperref}
\usepackage{parskip}
\usepackage{titlesec}
\usepackage{xcolor}
\usepackage{pdflscape}
\usepackage{natbib}

\geometry{margin=2.5cm}

\titleformat{\section}{\large\bfseries}{}{0em}{}[\titlerule]
\titleformat{\subsection}{\normalsize\bfseries}{}{0em}{}

\title{
    \textbf{Reporte de Avance}\\[0.5em]
    \large Masa del Gas Molecular en Envolventes Circumestelares:\\
    Tratamiento de la Opacidad y Revisión del Método
}
\author{Pablo Rojas}
\date{\today}

% ============================================================
\begin{document}
\maketitle
\tableofcontents
\newpage

% ============================================================
\section{Introducción y motivación}

Este reporte documenta la \textbf{revisión del cálculo de la masa del gas molecular}
($M_{\rm gas}$) en las envolventes circumestelares (CSE) de la muestra de nueve estrellas
AGB y post-AGB. Respecto a la versión anterior, el cambio central es el \textbf{tratamiento
físico de la opacidad} de la línea de CO: en lugar de adoptar un valor fijo y poco
justificado ($\tau = 1$), la opacidad se deriva ahora por fuente a partir de la razón
$^{12}$CO/$^{13}$CO observada.

El reporte se enfoca en tres puntos:
\begin{enumerate}
    \item \textbf{Cómo se calcula} la masa y cómo se determina la opacidad (Secciones~2 y 3).
    \item \textbf{Las complicaciones} encontradas —principalmente la elección de la abundancia
    isotópica $^{12}$CO/$^{13}$CO— y cómo se resolvieron (Sección~3).
    \item \textbf{Por qué el paper sugerido como referencia} (Xie, Young \& Schloerb 1994)
    \textbf{no es aplicable} a nuestro problema (Sección~4).
\end{enumerate}

% ============================================================
\section{Metodología del cálculo de masa}

\subsection{Fórmula general}

La masa de H$_2$ se obtiene a partir de la densidad de columna de CO integrada sobre
el haz \citep{choi1993, mangum2015}:

\begin{equation}
    M_{{\rm H}_2} = m_{{\rm H}_2}\, D^2\, \Omega\, \frac{N_{\rm CO}}{f_{\rm CO}},
    \qquad M_{\rm gas} = 1.36\,M_{{\rm H}_2},
    \label{eq:Mmain}
\end{equation}

\noindent donde $m_{{\rm H}_2}$ es la masa de la molécula de H$_2$, $D$ la distancia,
$\Omega$ el ángulo sólido del haz, $f_{\rm CO}=N_{\rm CO}/N_{\rm H_2}$ la abundancia de CO
(adoptada por tipo espectral; \citealp{ramstedt2008}), y el factor $1.36$ corrige por helio.

La densidad de columna se obtiene del flujo integrado de la línea, con la corrección de
opacidad explícita:

\begin{equation}
    N_{\rm CO} = \frac{C_J(T_{\rm ex})}{D_J(T_{\rm ex})}\, W_J \cdot
    \frac{\tau_J}{1 - e^{-\tau_J}},
    \qquad
    W_J = 1.222\times10^6\,\frac{I_{\rm CO}}{\nu^2\,\theta_{\rm maj}\,\theta_{\rm min}},
    \label{eq:NCO}
\end{equation}

\noindent donde $W_J$ es la temperatura de brillo integrada [K\,km\,s$^{-1}$], $I_{\rm CO}$ el
flujo integrado [Jy\,km\,s$^{-1}$], $C_J$ el prefactor de columna y $D_J$ la corrección de
Rayleigh-Jeans de la transición. El ángulo sólido del haz se cancela con $W_J$, por lo que
la masa de una fuente no resuelta no depende del tamaño del haz.

\subsection{Transición utilizada por fuente}

La opacidad $\tau$ pertenece a la transición en la que se mide la razón isotópica. Por ello,
\textbf{la masa se calcula desde la transición en la que se dispone de $^{13}$CO}:
\begin{itemize}
    \item \textbf{CO$(1\to0)$} para las fuentes con $^{13}$CO$(1\to0)$: R~Hor, RV~Aqr, RZ~Sgr, UU~For.
    \item \textbf{CO$(2\to1)$} para las que solo tienen $^{13}$CO$(2\to1)$: los cuatro IRAS y V510~Pup.
\end{itemize}
Esto mezcla transiciones entre fuentes, pero garantiza que $\tau$ y la masa sean
\emph{autoconsistentes} (medidas en la misma línea). Los prefactores correspondientes, a
$T_{\rm ex}=20$\,K, son:
\begin{equation}
    \frac{C_{10}}{D_{10}} = \frac{1.14\times10^{15}}{1.14} \approx 1.0\times10^{15},
    \qquad
    \frac{C_{21}}{D_{21}} = \frac{4.96\times10^{14}}{1.30} \approx 3.8\times10^{14}
    \;\;[\text{cm}^{-2}\,(\text{K\,km\,s}^{-1})^{-1}].
\end{equation}
El valor para $(1\to0)$ coincide con el $1.1\times10^{15}$ del paper del grupo para
$(3\to2)$, confirmando que la adaptación entre transiciones es correcta.

% ============================================================
\section{El tratamiento de la opacidad: complicaciones y discusión}

Esta sección es el núcleo del reporte, ya que la opacidad fue la principal dificultad.

\subsection{Por qué $\tau=1$ no era suficiente}

En la versión anterior se adoptó $\tau_{10}=1$ (corrección $\sim 1.58$) como valor
``moderado''. Esto es arbitrario: la corrección de opacidad va desde $1$ (delgado) hasta
$\gtrsim 3$ (grueso), y con $\tau=3$ la masa se duplica respecto a $\tau=1$. Es decir,
\textbf{la opacidad es uno de los sistemáticos dominantes} y no podía dejarse fija sin
justificación. A raíz de la discusión se decidió derivarla de los datos.

\subsection{Opacidad a partir de $^{12}$CO/$^{13}$CO y la relación $R \leq X$}

La opacidad se estima de la razón observada $R = T(^{12}{\rm CO})/T(^{13}{\rm CO})$
asumiendo que $^{13}$CO es ópticamente delgado y que ambas líneas comparten excitación
\citep{ramstedt2014}:
\begin{equation}
    R = \frac{1-e^{-\tau_{12}}}{1-e^{-\tau_{13}}},
    \qquad \tau_{13} = \tau_{12}/X,
    \label{eq:tau}
\end{equation}
donde $X = {}^{12}$CO/$^{13}$CO es la abundancia isotópica. De aquí se obtiene una propiedad
clave:
\begin{equation}
    R = X\cdot\frac{1-e^{-\tau_{12}}}{\tau_{12}} \;\leq\; X,
    \qquad\text{con } R = X \text{ en el límite delgado } (\tau_{12}\to0).
    \label{eq:RleqX}
\end{equation}
Es decir, \textbf{la razón observada nunca puede superar la abundancia real}: la opacidad
solo \emph{baja} $R$ por debajo de $X$. Una razón observada alta es, por tanto, la firma
directa de gas ópticamente delgado.

\subsection{La complicación principal: ¿qué valor de $X$ usar?}

El valor de $X$ no es observable directamente y debe adoptarse. La referencia natural es
\citet{ramstedt2014}, que da las medianas por tipo espectral: $X=13$ (M), $26$ (S) y $34$ (C).
\textbf{El problema, discutido en detalle, es que estas distribuciones son muy dispersas},
especialmente para las estrellas de carbono, donde el cociente entre los percentiles 90 y 10
es de un factor $\sim15$. Parte de esa dispersión proviene de las estrellas de tipo J
(con $X\approx3$--$4$), que son una clase distinta; excluyéndolas, las estrellas de carbono
normales (N-type) caen en $X\sim20$--$90$ \citep{lambert1986}. Una única mediana no es,
entonces, representativa de una estrella individual.

Esto produjo un caso aparentemente paradójico: IRAS\,06531 tiene $R=84.68$, mayor que la
mediana de carbono $X=34$, lo que por la Ec.~(\ref{eq:RleqX}) sería ``imposible''.

\subsection{Resolución: la regla \texorpdfstring{$X=\max(R,X_{\rm tipo})$}{X=max(R,Xtipo)} y la clasificación delgado/grueso}

La paradoja anterior \textbf{no indica una contradicción, sino que revela el valor real de
$X$ de esa estrella}: si $R\leq X$ siempre, entonces observar $R=84.68$ implica $X\geq84.68$.
La mediana $34$ simplemente subestima $X$ para esa fuente, que es una estrella de carbono de
$X$ alto (algo normal, hasta $\sim100$; \citealp{lambert1986}). Con $X\approx R$, la
Ec.~(\ref{eq:RleqX}) da $\tau\to0$: la fuente es \textbf{ópticamente delgada}. Es decir, una
razón alta no es un error, es la medición de un $X$ alto en gas delgado.

Esto se implementa como un \textbf{piso físico} sobre la abundancia:
\begin{equation}
    X_{\rm estrella} = \max\!\big(R_{\rm obs},\, X_{\rm tipo}\big).
    \label{eq:floor}
\end{equation}
La regla clasifica la opacidad directamente desde los datos:
\begin{itemize}
    \item $R$ alto (comparable o mayor que $X_{\rm tipo}$) $\Rightarrow$ \textbf{delgado} ($\tau\approx0$);
    la masa \emph{no depende} del valor exacto de $X$.
    \item $R$ bajo ($R \ll X_{\rm tipo}$) $\Rightarrow$ \textbf{grueso}; aquí $X$ sí importa.
\end{itemize}

\subsection{La masa como rango, no como un punto}

Para las fuentes gruesas, donde $X$ afecta el resultado, en lugar de un único valor se reporta
un \textbf{rango de masa} propagando el percentil de $X$ por tipo
(M: $[6.5,26]$; S: $[15,45]$; C: $[20,90]$), siempre con el piso de la Ec.~(\ref{eq:floor}).
Así, la dispersión de \citet{ramstedt2014} se convierte en una barra de error honesta en vez
de quedar oculta en una mediana. Para las fuentes delgadas el rango colapsa al valor central,
reflejando que ahí la incertidumbre en $X$ es irrelevante.

\subsection{Sensibilidad a la temperatura de excitación}

\citet{choi1993} advierte que la densidad de columna de CO$(1\to0)$ depende más de $T_{\rm ex}$
que la de $(3\to2)$. Verificamos que, en el rango realista $T_{\rm ex}=15$--$40$\,K, la masa
varía un factor $\sim2.4$ para $(1\to0)$ y $\sim1.6$ para $(2\to1)$, frente a $\sim1.2$ para
$(3\to2)$. Siguiendo a \citet{cerrigone2012} —que fija el factor de columna ``dentro de un
factor de 2''— adoptamos $T_{\rm ex}=20$\,K y reportamos esta variación como un sistemático de
$\sim$factor 2, sin intentar ajustar $T_{\rm ex}$.

% ============================================================
\section{Por qué el método de Xie et al.\ (1994) no es aplicable}

Se evaluó en detalle el trabajo de \citet{xie1994}, sugerido como referencia para el cálculo
de masa. Se concluye que \textbf{no es aplicable a nuestro problema}, por tres razones:

\begin{enumerate}
    \item \textbf{Es un estudio de galaxias} (NGC\,2146 e IC\,342), no de envolventes
    circumestelares. El objeto físico, las escalas y las condiciones del gas son distintos.

    \item \textbf{Su método de masa es el factor de conversión $N({\rm H_2})/I_{\rm CO}$}
    (de tipo $\alpha_{\rm CO}$ galáctico), calibrado para nubes moleculares del medio
    interestelar. Este es precisamente el método que \emph{descartamos} para CSE, donde no es
    válido. Nuestro cálculo usa densidad de columna por fuente (Choi/Cerrigone), que es lo
    apropiado para envolventes circumestelares.

    \item \textbf{Para la opacidad usa la misma razón $^{12}$CO/$^{13}$CO que nosotros, pero
    asume un valor fijo $X=60$} (valor galáctico). Es decir, \textbf{no resuelve} la dificultad
    central —cómo fijar $X$ sin un valor disperso y poco representativo—, y además el valor que
    usa no aplica a estrellas evolucionadas, donde $X$ depende fuertemente del tipo y es mucho
    menor.
\end{enumerate}

\noindent En resumen, el paper de \citet{xie1994} \textbf{confirma} que el diagnóstico correcto
de opacidad es la razón $^{12}$CO/$^{13}$CO (lo que ya hacíamos), pero \textbf{ni su método de
masa ni su valor de $X$ son trasladables} a nuestra muestra. La referencia metodológica
adecuada para masa con líneas bajas de CO en objetos evolucionados es \citet{huggins1996}
(que usa CO$(1\to0)$ y $(2\to1)$ en nebulosas planetarias) junto con \citet{mangum2015}.

% ============================================================
\section{Resultados}

La Tabla~\ref{tab:mass} presenta la masa de gas con su rango por opacidad y la clasificación
delgado/marginal/grueso resultante. El error $\sigma/M$ es el estadístico (flujo $+$ distancia);
el rango $[M_{\rm lo},M_{\rm hi}]$ refleja el sistemático por la abundancia $X$.

\begin{landscape}
\begin{table}[h]
\centering
\caption{Masa de gas molecular con tratamiento de opacidad. $R$ es la razón
$^{12}$CO/$^{13}$CO observada; $X_{\rm med}$ la abundancia central por tipo (con piso
$X\geq R$); $\tau$ la opacidad central; $M_{\rm gas}$ el valor central y
$[M_{\rm lo},M_{\rm hi}]$ el rango por el percentil de $X$. Régimen: T=delgado, M=marginal,
G=grueso. $T_{\rm ex}=20$\,K.}
\label{tab:mass}
\begin{tabular}{llcccccccl}
\toprule
Fuente & Trans. & Tipo & $D$ (pc) & $f_{\rm CO}$ & $R$ & $X_{\rm med}$ & $\tau$ &
$M_{\rm gas}$ ($M_\odot$) & $[M_{\rm lo},M_{\rm hi}]$ ($M_\odot$) / régimen \\
\midrule
IRAS 06531$-$0216 & 2$-$1 & C & 1158 & $1.0\!\times\!10^{-3}$ & 84.7 & 34 & 0.00 & $5.98\times10^{-4}$ & $[5.98,\,6.36]\times10^{-4}$ \;\;(T) \\
IRAS 07145$-$1428 & 2$-$1 & C & 4145 & $1.0\!\times\!10^{-3}$ & 26.7 & 34 & 0.52 & $4.28\times10^{-3}$ & $[3.33,\,11.4]\times10^{-3}$ \;\;(M) \\
IRAS 07180$-$1314 & 2$-$1 & M & 1168 & $1.0\!\times\!10^{-3}$ & 4.95 & 13 & 2.72 & $2.28\times10^{-4}$ & $[1.09,\,4.60]\times10^{-4}$ \;\;(G) \\
IRAS 08500$-$3254 & 2$-$1 & M & 1527 & $2.0\!\times\!10^{-4}$ & 4.42 & 13 & 3.18 & $3.11\times10^{-2}$ & $[1.49,\,6.25]\times10^{-2}$ \;\;(G) \\
R Hor             & 1$-$0 & M &  231 & $2.0\!\times\!10^{-4}$ & 39.1 & 13 & 0.00 & $1.00\times10^{-3}$ & punto \;\;(T) \\
RV Aqr            & 1$-$0 & C &  635 & $2.0\!\times\!10^{-4}$ & 27.1 & 34 & 0.49 & $2.68\times10^{-2}$ & $[2.12,\,7.18]\times10^{-2}$ \;\;(M) \\
RZ Sgr            & 1$-$0 & S &  424 & $6.0\!\times\!10^{-4}$ & 4.67 & 26 & 6.25 & $7.03\times10^{-3}$ & $[4.04,\,12.2]\times10^{-3}$ \;\;(G) \\
UU For            & 1$-$0 & M &  615 & $2.0\!\times\!10^{-4}$ & 18.8 & 13 & 0.00 & $6.41\times10^{-3}$ & $[6.41,\,8.98]\times10^{-3}$ \;\;(T) \\
V510 Pup          & 2$-$1 & post-AGB & 436 & $4.0\!\times\!10^{-4}$ & 12.1 & --- & 1.00$^{\dagger}$ & $6.65\times10^{-4}$ & provisional \\
\bottomrule
\end{tabular}
\\[0.3em]
{\footnotesize $^{\dagger}$ V510 Pup (post-AGB) no tiene $^{12}$CO/$^{13}$CO publicado
(fuera de la muestra de \citealp{ramstedt2014}); se adopta $\tau=1$ provisional. Su razón
observada ($12.1$) es un límite inferior de $X$, sugiriendo gas casi delgado.}
\end{table}
\end{landscape}

\noindent De las nueve fuentes: \textbf{3 son ópticamente delgadas} (IRAS\,06531, R\,Hor,
UU\,For), donde la masa no depende de $X$; \textbf{2 son marginales} (IRAS\,07145, RV\,Aqr),
con un rango amplio que abarca desde delgado hasta moderadamente grueso; \textbf{3 son gruesas}
(IRAS\,07180, IRAS\,08500, RZ\,Sgr), donde la corrección de opacidad aumenta la masa de forma
significativa (hasta $\sim6\times$ en RZ\,Sgr); y V510\,Pup queda provisional.

% ============================================================
\section{Discusión y limitaciones}

\textbf{Qué cambió respecto a $\tau=1$.} Las fuentes delgadas tienen ahora masas \emph{menores}
(el $\tau=1$ anterior sobreestimaba), mientras que las gruesas tienen masas \emph{mayores} y
mejor justificadas. Sobre todo, la opacidad ya no es un número impuesto sino una cantidad
derivada de los datos.

\textbf{Limitación principal.} Para las fuentes gruesas, $\tau$ (y por tanto la masa) depende
$\sim$linealmente de la abundancia $X$, que se adopta por tipo. Lo robusto es la
\emph{clasificación} delgado/grueso, que sale directamente de $R$; el \emph{valor} de la masa
de las gruesas hereda la incertidumbre de $X$, que reportamos como rango. La forma rigurosa de
eliminar esta dependencia sería un modelo de transferencia radiativa completo de la envolvente
por fuente (estilo \citealp{ramstedt2014}), que derivaría $X$ de los datos; un código LVG
simple (p.\,ej.\ RADEX) no basta, porque también requiere asumir $X$ para relacionar las
columnas de $^{12}$CO y $^{13}$CO.

\textbf{Próximos pasos.} (i) Buscar valores de $X$ publicados por fuente (fotosféricos o de
transferencia radiativa) para las tres gruesas, donde $X$ es determinante; (ii) refinar
$\sigma_I$ con el ruido por canal de los cubos ALMA; (iii) obtener un $^{12}$CO/$^{13}$CO para
V510\,Pup.

% ============================================================
\begin{thebibliography}{9}

\bibitem[Cerrigone et al.(2012)]{cerrigone2012}
Cerrigone, L., Menten, K.~M., \& Kami\'nski, T.\
2012, \textit{A\&A}, 542, A15

\bibitem[Choi et al.(1993)]{choi1993}
Choi, M., Evans, N.~J.~II, \& Jaffe, D.~T.\
1993, \textit{ApJ}, 417, 241

\bibitem[Huggins et al.(1996)]{huggins1996}
Huggins, P.~J., Bachiller, R., Cox, P., \& Forveille, T.\
1996, \textit{A\&A}, 315, 284

\bibitem[Lambert et al.(1986)]{lambert1986}
Lambert, D.~L., Gustafsson, B., Eriksson, K., \& Hinkle, K.~H.\
1986, \textit{ApJS}, 62, 373

\bibitem[Mangum \& Shirley(2015)]{mangum2015}
Mangum, J.~G., \& Shirley, Y.~L.\
2015, \textit{PASP}, 127, 266

\bibitem[Ramstedt et al.(2008)]{ramstedt2008}
Ramstedt, S., Sch\"oier, F.~L., Olofsson, H., \& Lundgren, A.~A.\
2008, \textit{A\&A}, 487, 645

\bibitem[Ramstedt \& Olofsson(2014)]{ramstedt2014}
Ramstedt, S., \& Olofsson, H.\
2014, \textit{A\&A}, 566, A145

\bibitem[Xie et al.(1994)]{xie1994}
Xie, S., Young, J.~S., \& Schloerb, F.~P.\
1994, \textit{ApJ}, 421, 434

\end{thebibliography}

\end{document}
